% REVISED: canonical-target reframe 2026-06-04
% REVISED: hostile-review action plan 2026-06-04
%
\section{Introduction}
\label{sec:intro}

Cartel enforcement begins before legal proof exists. An
agency deciding where to investigate rarely starts with the
communications, bid histories, leniency testimony, or
tender-level conduct evidence that would support liability.
It starts with thinner information: administrative records,
complaints, procurement anomalies, repeated interactions, and
case leads that make some firms and markets worth a closer
look. The first legal-economic problem is therefore not proof
itself. It is the allocation of proof-producing effort. The
Becker--Stigler tradition treats investigation and sanctions as
costly instruments whose intensity must be chosen optimally
\citep{becker1968crime,stigler1970optimum}; the margin we study
is the one prior to any sanction---how a resource-constrained
authority sequences costly evidence acquisition before it knows
whom to charge. This is the enforcement-side counterpart to the
procurement-governance questions in \citet{decarolis2020bureaucratic},
where bureaucratic competence and discretion shape award
outcomes: there the question is how a buyer runs a tender well;
here it is how an enforcer decides where to spend forensic
capacity once the tenders are run.

That distinction matters in procurement. The empirical
cartel-screening literature usually starts from the bid layer:
submitted bid amounts, ranks, dispersion, histories, and
sometimes bidder networks. Classic screens compare bidding
under competition and collusion
\citep{porter1993detection,bajari2003deciding}; variance and
descriptive screens formalize suspicious bid distributions
\citep{abrantesmetz2006a,harrington2008detecting} and
field-test them in procurement
\citep{imhof2018screening,imhof2019detecting}; recent work
extends the approach through bidder groups
\citep{conley2016detecting}, machine learning
\citep{huber2019machine,wallimann2023machine}, and robust
screen design \citep{chassang2022robust}. Those objects are
close to the evidentiary stage but not always cheap at scale.
This is the central asymmetry the paper exploits: many
procurement systems make contract-award records routinely
visible---who participated, who won, item codes, and negotiated
prices---while keeping per-bidder bid files in operational
systems that must be recovered, cleaned, and inspected case by
case. The award layer is cheap and broad but legally thin; the
bid layer is forensically rich but costly to open.

Cheap administrative red flags built from exactly such routine
records are already in institutional use. Corruption-risk
indicators such as single bidding and repeat-participation
anomalies are computed from contract-award data across entire
procurement systems and used by multilateral institutions and
audit bodies precisely because bid-level records are costly to
recover \citep{fazekas2020uncovering,oecd2022datascreening}. The
zero-win construct we study is the cartel-facing member of this
cheap-red-flag family: a statistic an authority can run over
award data before committing to forensic recovery. The paper
audits the reach of this family before it steers scarce forensic
effort. An enforcement agency thus faces a sequencing
problem---before opening the costly bid layer, can it use the
cheap award layer to decide where bid-layer forensics should
begin?

This paper studies that question on S\~ao Paulo's BEC
procurement platform, covering $\valSampleN$ tender-items from
$\valYearStart$ to $\valYearEnd$. BEC records a routine
award layer---participants, winners, item codes, negotiated
prices, procuring units, years, and modalities---and also
preserves a richer bid layer recoverable through the LANCES
export. CADE procurement-cartel adjudications provide external
legal anchors. This combination lets us study an earlier stage
than most bid-distribution screens: the point at which an agency
observes award records but has not yet recovered the bid-level
evidence needed for forensic evaluation.

The award-layer statistic is deliberately simple. Among firms
with zero wins, we rank participation intensity using
$\log(1+\text{tenders\_count})$ and implement an auditable
binary rule that flags \emph{frequent losers}. The logic is
monotone rather than classificatory: if coordinated bidders use
loser-side participation to create the appearance of competition
while avoiding the win, persistent zero-win participation can
rank loser-side exposure under a transparent behavioral
condition. The statistic does not identify a firm type or assign
cartel membership; the rank allocates forensic attention, while
the richer bid layer remains where tender-level conduct and
agreement are evaluated.

That legal scope determines the validation target. Direct CADE
defendants are legal anchors, but they are not the natural object
of a zero-win loser-side screen. Procurement cartels often
allocate winning roles, rotate designated winners, and use other
firms on the losing side of the tender
\citep{pesendorfer2000study,asker2010leniency}; the same
role-allocation logic is central in cartel theory and case
evidence \citep{marshall2012economics}, and large-scale
procurement evidence shows how such arrangements appear in
administrative bidding records \citep{kawai2022detecting}. We
therefore validate the ranking against reproducible
adjudication-anchored exposure rather than against cartel
membership: the target is the set of $\valMainCobidders$ unique
always-loser firms that share at least one tender-item with a
BEC-active direct CADE defendant, with the defendants themselves
excluded. The label is rebuilt end-to-end from the current
pipeline scripts, and---this point matters for the validity of
every comparison that follows---the frequent-loser flag never
enters its construction. Direct defendants anchor legally
adjudicated environments; the always-loser cobidders are the
loser-side footprint the award layer can plausibly rank before
bid-level proof is recovered.

The empirical design is organized around three threats that would
each, on its own, make such a ranking uninformative. The first is
opportunity arithmetic: high-volume zero-win firms mechanically
meet more defendants, so a raw co-bidding association may reflect
participation volume and opportunity-set exposure rather than any
behavioral signal. The second is timing: a ranking formed using
contemporaneous or later participation cannot order forensic
priority prospectively. The third is case concentration: a single
adjudicated case can manufacture apparent platform-wide reach. We
discipline the ranking against all three---opportunity-cell
adjustment, strict rolling-origin timing, and
leave-one-case-out tests---alongside a bid-layer benchmark;
Section~4 and the appendices report the battery, and
\ref{app:validation_audits_submission} gives the exact exposure,
leakage, and timing audit construction.

The results are deflationary, and that is the point. Read raw,
the award-layer ranking does concentrate adjudication-anchored
cobidders inside the always-loser stratum; the decomposition then
shows how little survives discipline. Most of the raw
concentration reflects procurement opportunity, because firms that
participate more, in the same environments as adjudicated
defendants, are mechanically more exposed to them. Within
comparable opportunity sets the residual ordering is marginal at
best and not robust across designs: held to matched cells the
ranking is close to chance, the increment over an exposure-only
model is small, and matched-stratum permutation and enrichment
tests do not reject the opportunity-only null. Strict prospective
ranking fails outside the incumbent pool, and the apparent reach
is case-sensitive---a single adjudicated case supplies a large
share of the positives and of the firms a fixed-budget queue would
open. The ranking is also silent about direct CADE defendants,
consistent with its loser-side scope rather than any
cartel-membership reading.

%
Where the framework earns its keep is in cost accounting and
forensic sequencing. On the same adjudication-anchored target the
award-layer ranking is comparable to a transparent bid-moment
random-forest benchmark inspired by Imhof--Wallimann-style
screens, and the two combine only conditionally---a combined model
improves precision under random cross-validation but falls below
the award-only ranking under case-grouped folds. Used as a
gatekeeper, the award layer traces a cost-recall frontier: across
survivor-pool sizes it retains much of the adjudication-anchored
recall of full bid-layer scoring while shrinking the bid-recovery
footprint. The reduction is real but must be read against its
denominator---large in firms opened, far smaller in tender-items
and bid rows, because the survivors are high-participation
firms---and there is no optimal cutoff: the frontier itself, not
any single operating point, is the enforcement-design object.
Because strict timing for the bid rerank is not available with the
current bid-layer features, this is a retrospective cost-footprint
design conditional on the validated incumbent ranking, not a fully
prospective deployment test.

This paper makes three contributions. The first is
organizational, and it is the law-economics-organization payload:
the award-to-bid recovery decision is a sequential
evidence-allocation problem, and the object an enforcement
designer chooses is not any cutoff on the cheap screen but the
cost-recall frontier that the screen traces when used as a
gatekeeper. Reframing procurement-cartel screening this way turns
it from a classification exercise---does the flag attain a high
AUC---into a sequencing problem: how a resource-constrained
authority orders costly evidence acquisition under incomplete
observability, where suspicion is cheap and proof is dear
\citep{becker1968crime,stigler1970optimum,harrington2008detecting,sanchezgraells2019screening,chassang2022robust}.
Existing bid-distribution screens evaluate suspicious bidding once
bid-level data have been recovered; we study the prior stage,
where agencies observe award records but must decide where
recovery should start, and we show that the frontier---not a
detector-versus-detector horse race---is the design margin.

The second contribution is a disciplined audit protocol whose
components are individually standard but whose sequencing for the
award-to-bid decision is, to our knowledge, new: a non-circular
adjudication-anchored label, opportunity-cell adjustment, strict
rolling-origin timing, leave-one-case-out tests, a bid-layer
benchmark, and cost-recall accounting, applied to ask not whether
a cheap screen attains a high pooled AUC but how much of that AUC
is genuine ranking signal rather than mechanical opportunity
exposure, retrospective information, or concentration in a single
adjudicated case. We do not claim a validated, transferable
\emph{score}; what is portable is the \emph{protocol}---validating
an administrative screen against adjudicated cases without
adjusting for procurement opportunity systematically over-credits
it---and the order in which the audit steps must be run. We
demonstrate that portability directly: re-running the identical
battery on a second platform, the federal ComprasNet system
(Section~\ref{sec:comparative}), with partially overlapping CADE
anchors and no bid microdata, returns the same deflationary
verdict, so it is the audit, not the raw screen, that travels.

Third, applying that protocol to S\~ao Paulo's BEC yields a map
of reach and limits for this setting: persistent zero-win
participation concentrates a legally limited,
adjudication-anchored exposure target in raw terms, but the
residual ordering net of opportunity is marginal at best,
prospective ranking fails outside incumbent pools, and the
apparent reach is concentrated in one case---so the estimated
ranking is not a portable cartel score, while the audit principle
and the cost-recall trade-offs it exposes are what transfer.
Locating these boundaries is itself the enforcement-design
product.

Two scope points round out the design. Price evidence enters only
as scope information, not as a damages exercise. And strategic
response shapes implementation: a static cutoff invites gaming,
while refreshed thresholds, continuous ranks or
capacity-constrained queues, bunching checks, and bid-layer
follow-up fit the staged enforcement architecture.

The broader lesson is institutional. A legal system that waits
for full proof before ranking investigative priorities will
spend scarce forensic capacity poorly; a system that treats cheap
suspicion as proof will overreach. The useful middle ground is a
disciplined triage rule---cheap enough to run before proof
exists, informative enough to order forensic attention within a
known reach, and legally limited enough to leave liability where
it belongs, in the richer evidentiary record. Knowing where such
a rule stops ranking is as much a part of disciplined enforcement
as knowing where it begins.
